\begin{theorem}[单元曲率的二矩完备性]\label{thm:conclusion-cell-curvature-two-moment-completeness}
\leanverified{paper\_conclusion\_cell\_curvature\_two\_moment\_completeness}
设 \(n\ge 1\) 为整数，并假设
\[
\operatorname{supp}\mu_\Lambda\subset [n,n+1].
\]
则恰有两条离散扇宽可能非零，即 \(w_{n+1}\) 与 \(w_{n+2}\)，并且它们精确给出该单元内曲率测度的总质量与一阶矩：
\[
w_{n+1}=\int_{[n,n+1]}(n+1-t)\,d\mu_\Lambda(t),
\qquad
w_{n+2}=\int_{[n,n+1]}(t-n)\,d\mu_\Lambda(t).
\]
从而
\[
w_{n+1}+w_{n+2}=\mu_\Lambda([n,n+1]),
\]
以及
\[
n\,w_{n+1}+(n+1)\,w_{n+2}
=
\int_{[n,n+1]} t\,d\mu_\Lambda(t).
\]
\end{theorem}
\begin{proof}
由定理 \ref{thm:conclusion-discrete-fan-width-tent-kernel-tomography}，
\[
w_q=\int_\Theta K\bigl(t-(q-1)\bigr)\,d\mu_\Lambda(t),
\qquad
K(u)=(1-|u|)_+.
\]
由于 \(K(\cdot-(q-1))\) 的支撑为 \([q-2,q]\)，若 \(\operatorname{supp}\mu_\Lambda\subset[n,n+1]\)，则只有 \(q=n+1,n+2\) 的帐篷核与该单元相交。并且在 \([n,n+1]\) 上，
\[
K(t-n)=n+1-t,
\qquad
K\bigl(t-(n+1)\bigr)=t-n.
\]
代入即可得到前两式。

两式相加给出
\[
w_{n+1}+w_{n+2}
=
\int_{[n,n+1]}\bigl((n+1-t)+(t-n)\bigr)\,d\mu_\Lambda(t)
=
\mu_\Lambda([n,n+1]).
\]
再作线性组合
\[
n(n+1-t)+(n+1)(t-n)=t,
\]
对 \(\mu_\Lambda\) 积分即得一阶矩公式。证毕。
\end{proof}

\begin{corollary}[单元内曲率重心的邻峰比值公式]\label{cor:conclusion-cell-curvature-barycenter-adjacent-peak-ratio}
\leanverified{paper\_conclusion\_cell\_curvature\_barycenter\_adjacent\_peak\_ratio}
在定理 \ref{thm:conclusion-cell-curvature-two-moment-completeness} 的设定下，若
\[
\mu_\Lambda([n,n+1])>0,
\]
则该单元内曲率测度的重心
\[
\bar t_n:=
\frac{\int_{[n,n+1]} t\,d\mu_\Lambda(t)}
{\mu_\Lambda([n,n+1])}
\]
可由离散扇宽精确恢复为
\[
\bar t_n
=
n+\frac{w_{n+2}}{w_{n+1}+w_{n+2}}.
\]
\end{corollary}
\begin{proof}
由定理 \ref{thm:conclusion-cell-curvature-two-moment-completeness}，
\[
\int_{[n,n+1]} t\,d\mu_\Lambda(t)
=
n\,w_{n+1}+(n+1)\,w_{n+2},
\]
且
\[
\mu_\Lambda([n,n+1])=w_{n+1}+w_{n+2}.
\]
两式相除即可。证毕。
\end{proof}

\begin{theorem}[单隐藏相变点的精确反演]\label{thm:conclusion-single-hidden-phase-transition-exact-inversion}
\leanverified{paper\_conclusion\_single\_hidden\_phase\_transition\_exact\_inversion}
沿用定理 \ref{thm:conclusion-cell-curvature-two-moment-completeness} 的设定。若更强地
\[
\mu_\Lambda=a\,\delta_\tau,
\qquad
\tau\in[n,n+1],\quad a>0,
\]
则离散扇宽满足
\[
w_{n+1}=a(n+1-\tau),
\qquad
w_{n+2}=a(\tau-n).
\]
因而该隐藏相变点的强度与位置可显式反演为
\[
a=w_{n+1}+w_{n+2},
\qquad
\tau=n+\frac{w_{n+2}}{w_{n+1}+w_{n+2}}.
\]
\end{theorem}
\begin{proof}
把 \(\mu_\Lambda=a\delta_\tau\) 代入定理 \ref{thm:conclusion-cell-curvature-two-moment-completeness} 的两条公式，立即得到
\[
w_{n+1}=a(n+1-\tau),
\qquad
w_{n+2}=a(\tau-n).
\]
两式相加给出 \(a\)；再由第二式除以和式得到 \(\tau\)。证毕。
\end{proof}

\begin{theorem}[尾曲率预算与连续冻结的等价]\label{thm:conclusion-tail-curvature-budget-equivalent-continuous-freezing}
对任意整数 \(N\ge 2\)，定义尾部权函数
\[
\eta_N(t):=\sum_{q=N}^{\infty}K\bigl(t-(q-1)\bigr).
\]
则有精确闭式
\[
\eta_N(t)=
\begin{cases}
0, & t\le N-2,\\[2mm]
t-(N-2), & N-2\le t\le N-1,\\[2mm]
1, & t\ge N-1.
\end{cases}
\]
从而
\[
\sum_{q=N}^{\infty}w_q
=
\int_{N-2}^{N-1}\bigl(t-(N-2)\bigr)\,d\mu_\Lambda(t)
\;+\;
\mu_\Lambda([N-1,\infty)).
\]
特别地，成立严格等价：
\[
w_q=0\quad(\forall q\ge N)
\iff
\mu_\Lambda\bigl((N-2,\infty)\bigr)=0.
\]
等价地，\(\Lambda\) 在开半轴 \((N-2,\infty)\) 上恰为仿射函数。若再有 \(\mu_\Lambda(\{N-2\})=0\)，则 \(\Lambda\) 在闭半轴 \([N-2,\infty)\) 上仿射。
\leanverified{paper\_conclusion\_tail\_curvature\_budget\_equivalent\_continuous\_freezing}
\end{theorem}
\begin{proof}
对 \(t\le N-2\)，所有中心 \(q-1\ge N-1\) 与 \(t\) 的距离都至少为 \(1\)，故 \(\eta_N(t)=0\)。

若 \(t\in[N-2,N-1]\)，则只有 \(q=N\) 这一项非零，于是
\[
\eta_N(t)=K\bigl(t-(N-1)\bigr)=t-(N-2).
\]

若 \(t\ge N-1\)，则所有可能非零的帐篷核中心都已落在求和范围之内；由平移帐篷核的分片线性分割恒等式，
\[
\sum_{m\in\ZZ}K(t-m)=1,
\]
且对给定 \(t\ge N-1\) 的非零项都满足 \(m\ge N-1\)，故 \(\eta_N(t)=1\)。

再由 Tonelli 定理交换求和与积分，
\[
\sum_{q=N}^{\infty}w_q
=
\int \eta_N(t)\,d\mu_\Lambda(t),
\]
代入上述闭式便得尾预算公式。

由于 \(\eta_N(t)\ge 0\)，且在 \((N-2,\infty)\) 上严格为正，故
\[
\sum_{q=N}^{\infty}w_q=0
\iff
\mu_\Lambda\bigl((N-2,\infty)\bigr)=0.
\]
又因 \(w_q\ge 0\)，和为零当且仅当每一项都为零，即
\[
\sum_{q=N}^{\infty}w_q=0
\iff
w_q=0\quad(\forall q\ge N).
\]
最后，凸函数在某开区间上二阶分布导数为零当且仅当其在该区间上仿射；若端点 \(N-2\) 也无原子，则仿射性可延伸到闭半轴。证毕。
\end{proof}

\begin{theorem}[孤立扇峰的整数刚性]\label{thm:conclusion-isolated-fan-peak-integer-rigidity}
\leanverified{paper\_conclusion\_isolated\_fan\_peak\_integer\_rigidity}
设 \(q\ge 3\)。若出现一个孤立非零扇峰
\[
w_{q-1}=0,\qquad
w_q>0,\qquad
w_{q+1}=0,
\]
则在局部窗口 \([q-2,q]\) 内，曲率测度被强制坍缩为单点原子：
\[
\mu_\Lambda\bigl([q-2,q]\setminus\{q-1\}\bigr)=0,
\qquad
\mu_\Lambda(\{q-1\})=w_q.
\]
因此 \(t=q-1\) 是一个精确整数相变点，其跳跃质量正好等于该孤立扇峰的高度。
\end{theorem}
\begin{proof}
对 \(t\in[q-2,q-1)\)，有
\[
K\bigl(t-(q-2)\bigr)=q-1-t>0.
\]
若该区间上有正质量，则
\[
w_{q-1}=\int K\bigl(t-(q-2)\bigr)\,d\mu_\Lambda(t)>0,
\]
与 \(w_{q-1}=0\) 矛盾。因此
\[
\mu_\Lambda([q-2,q-1))=0.
\]
同理，对 \(t\in(q-1,q]\)，有
\[
K(t-q)=t-q+1>0.
\]
若该区间上有正质量，则
\[
w_{q+1}=\int K(t-q)\,d\mu_\Lambda(t)>0,
\]
与 \(w_{q+1}=0\) 矛盾。故
\[
\mu_\Lambda((q-1,q])=0.
\]
于是 \([q-2,q]\) 内唯一可能的支撑点是 \(q-1\)。因此
\[
w_q=\int K\bigl(t-(q-1)\bigr)\,d\mu_\Lambda(t)=K(0)\mu_\Lambda(\{q-1\})=\mu_\Lambda(\{q-1\}),
\]
因为 \(K(0)=1\)。证毕。
\end{proof}

\begin{theorem}[整数支撑曲率的精确样条重建]\label{thm:conclusion-integer-supported-curvature-exact-spline-reconstruction}
\leanverified{paper\_conclusion\_integer\_supported\_curvature\_exact\_spline\_reconstruction}
假设
\[
\mu_\Lambda((0,1))=0,
\qquad
\operatorname{supp}\mu_\Lambda\subset \ZZ_{\ge 1}.
\]
则存在一列非负系数 \((\kappa_n)_{n\ge 1}\) 使
\[
\mu_\Lambda=\sum_{n\ge 1}\kappa_n\delta_n,
\qquad
\kappa_n=w_{n+1}.
\]
并且整个连续压力曲线可由离散扇宽显式重建为
\[
\Lambda(t)
=
\Lambda(0)+\bigl(\Lambda(1)-\Lambda(0)\bigr)t
\;+\;
\sum_{n\ge 1}w_{n+1}(t-n)_+,
\qquad
t\ge 0.
\]
因此，在整数支撑制度下，离散扇宽序列 \((w_q)_{q\ge 2}\) 与两个端点锚 \(\Lambda(0),\Lambda(1)\) 已足以唯一决定整条连续压力曲线与其 Legendre 前沿。
\end{theorem}
\begin{proof}
由支撑假设，可写
\[
\mu_\Lambda=\sum_{n\ge 1}\kappa_n\delta_n.
\]
于是对每个 \(q\ge 2\)，
\[
w_q
=
\sum_{n\ge 1}\kappa_n K\bigl(n-(q-1)\bigr).
\]
而当 \(n,q\) 为整数时，
\[
K\bigl(n-(q-1)\bigr)=
\begin{cases}
1,& n=q-1,\\
0,& n\neq q-1,
\end{cases}
\]
故
\[
w_q=\kappa_{q-1},
\qquad
\kappa_n=w_{n+1}.
\]

再定义
\[
\Psi(t):=
\Lambda(0)+\bigl(\Lambda(1)-\Lambda(0)\bigr)t
\;+\;
\sum_{n\ge 1}w_{n+1}(t-n)_+.
\]
对每个固定 \(t\)，上式只有有限项非零，故 \(\Psi\) 良定义且为凸分段仿射函数。其分布二阶导数满足
\[
\Psi''=\sum_{n\ge 1}w_{n+1}\delta_n=\mu_\Lambda=\Lambda''.
\]
因此 \(\Lambda-\Psi\) 的二阶分布导数为零，故 \(\Lambda-\Psi\) 为仿射函数。又因 \(\mu_\Lambda((0,1))=0\)，\(\Lambda\) 在 \([0,1]\) 上仿射；而 \(\Psi(0)=\Lambda(0)\)、\(\Psi(1)=\Lambda(1)\)，故这条仿射差在两个点上为零，只能恒为零。于是 \(\Lambda\equiv\Psi\)。证毕。
\end{proof}

\begin{corollary}[孤立峰族的全局 polyhedral spline 闭包]\label{cor:conclusion-isolated-peak-family-global-polyhedral-spline-closure}
\leanverified{paper\_conclusion\_isolated\_peak\_family\_global\_polyhedral\_spline\_closure}
若除边界胞 \((0,1)\) 外，所有正扇峰都满足定理 \ref{thm:conclusion-isolated-fan-peak-integer-rigidity} 的孤立条件，则
\[
\operatorname{supp}\mu_\Lambda\cap[1,\infty)\subset \ZZ_{\ge 1},
\]
从而定理 \ref{thm:conclusion-integer-supported-curvature-exact-spline-reconstruction} 的整数支撑假设在尾部自动成立。于是 \(\Lambda\) 在 \([1,\infty)\) 上成为一个由离散扇峰完全决定的 polyhedral spline。
\end{corollary}
\begin{proof}
若某个整数 \(n\ge 1\) 不是曲率支撑点，则对应 \(w_{n+1}=0\)。若 \(w_{n+1}>0\)，按假设其为孤立扇峰，定理 \ref{thm:conclusion-isolated-fan-peak-integer-rigidity} 蕴含
\[
\mu_\Lambda([n-1,n+1]\setminus\{n\})=0.
\]
因此每个正扇峰只对应一个整数原子点，且尾部不存在非整数支撑。再由定理 \ref{thm:conclusion-integer-supported-curvature-exact-spline-reconstruction} 即得结论。证毕。
\end{proof}

\begin{theorem}[近孤立扇峰的定量局部化]\label{thm:conclusion-near-isolated-fan-peak-quantitative-localization}
\leanverified{paper\_conclusion\_near\_isolated\_fan\_peak\_quantitative\_localization}
设 \(q\ge 3\)、\(\varepsilon\in(0,1)\)。定义局部离心质量
\[
M^-_{q,\varepsilon}:=
\mu_\Lambda\bigl([q-2,\ q-1-\varepsilon]\bigr),
\qquad
M^+_{q,\varepsilon}:=
\mu_\Lambda\bigl([q-1+\varepsilon,\ q]\bigr).
\]
则有精确上界
\[
M^-_{q,\varepsilon}\le \frac{w_{q-1}}{\varepsilon},
\qquad
M^+_{q,\varepsilon}\le \frac{w_{q+1}}{\varepsilon},
\]
从而
\[
\mu_\Lambda\Bigl([q-2,q]\setminus(q-1-\varepsilon,\ q-1+\varepsilon)\Bigr)
\le
\frac{w_{q-1}+w_{q+1}}{\varepsilon}.
\]
\end{theorem}
\begin{proof}
对 \(t\in[q-2,q-1-\varepsilon]\)，有
\[
K\bigl(t-(q-2)\bigr)=q-1-t\ge \varepsilon.
\]
因此
\[
w_{q-1}
=
\int K\bigl(t-(q-2)\bigr)\,d\mu_\Lambda(t)
\ge
\int_{[q-2,q-1-\varepsilon]}\varepsilon\,d\mu_\Lambda(t)
=
\varepsilon M^-_{q,\varepsilon},
\]
故
\[
M^-_{q,\varepsilon}\le \frac{w_{q-1}}{\varepsilon}.
\]
同理，对 \(t\in[q-1+\varepsilon,q]\)，有
\[
K(t-q)=t-q+1\ge \varepsilon,
\]
于是
\[
w_{q+1}\ge \varepsilon M^+_{q,\varepsilon},
\qquad
M^+_{q,\varepsilon}\le \frac{w_{q+1}}{\varepsilon}.
\]
两式相加即得总离心质量上界。证毕。
\end{proof}

\begin{corollary}[近孤立扇峰的弱骨架化]\label{cor:conclusion-near-isolated-fan-peak-weak-skeletonization}
\leanverified{paper\_conclusion\_near\_isolated\_fan\_peak\_weak\_skeletonization}
设 \(q_j\to\infty\) 为一列整数，并记
\[
\nu_j:=
\frac{\mu_\Lambda|_{[q_j-2,q_j]}}{\mu_\Lambda([q_j-2,q_j])},
\qquad
\mu_\Lambda([q_j-2,q_j])>0.
\]
若
\[
\frac{w_{q_j-1}+w_{q_j+1}}{\mu_\Lambda([q_j-2,q_j])}\longrightarrow 0,
\]
则平移后的局部曲率分布
\[
E\longmapsto \nu_j(E+q_j-1)
\]
弱收敛到 \(\delta_0\)。换言之，局部曲率分布骨架化到整数中心 \(q_j-1\)。
\end{corollary}
\begin{proof}
由定理 \ref{thm:conclusion-near-isolated-fan-peak-quantitative-localization}，对任意固定 \(\varepsilon\in(0,1)\)，
\[
\nu_j\Bigl([q_j-2,q_j]\setminus(q_j-1-\varepsilon,q_j-1+\varepsilon)\Bigr)
\le
\frac{w_{q_j-1}+w_{q_j+1}}{\varepsilon\,\mu_\Lambda([q_j-2,q_j])}
\longrightarrow 0.
\]
因此任意 \(\varepsilon\)-邻域外的质量都趋于零，故 \(\nu_j\) 弱收敛到 \(\delta_{q_j-1}\)。平移后即得到 \(\delta_0\)。证毕。
\end{proof}
